% Six substantive propositions and one novelty corollary.
\section{Required Baseline Propositions}
\label{sec:p08-propositions}

This section states exactly six substantive propositions. The novelty result
is a corollary of the MAH-relevant-project-set identity, not a seventh
substantive proposition. All policy comparisons are finite comparisons between
\(M=0\) and \(M=1\); the binary regime is never differentiated.

\subsection{Proposition 1: Organizational sorting}
\label{prop:p08-sorting}

\paragraph{Assumptions and fixed objects.}
Fix a project \((q,m)\), a finite institutional wedge, regime \(M=1\), and a
conjectured support price \(p_m<\infty\). Hold all primitives except internal
manufacturing capability \(k_i\) fixed. On the internally feasible domain,
use the technology signs and endpoint-crossing conditions stated above.

\paragraph{Claim.}
The internal-minus-entrusted value gap is strictly increasing in \(k_i\):
\begin{equation}
  \frac{\partial\Delta_{IE}(k_i;q,m,M,p_m)}{\partial k_i}
  =
  s(q)R_c(q,c_I)c_{I,k}(m,k_i)-F_{I,k}(m,k_i)
  >0.
  \label{eq:p08-sorting-slope}
\end{equation}
If the gap is negative at the lower feasible endpoint and positive at the
upper endpoint, there is a unique \(k^*(q,m;p_m,M)\). Conditional on both
\(I\) and \(E\) dominating \(T\) and \(A\), developers below the cutoff use
\(E\) and developers above it use \(I\).

\paragraph{Proof.}
The commercial-value kernel gives \(R_c<0\), while the internal technology
gives \(c_{I,k}<0\) and \(F_{I,k}<0\).
The operating-cost term on the right side of
\eqref{eq:p08-sorting-slope} is weakly positive (and is zero when
\(s(q)=0\)), while the setup-cost term is strictly positive. Their sum is
therefore strictly positive. Continuity and endpoint crossing give existence
by the intermediate value theorem; strict monotonicity gives uniqueness. The sign of
the gap on either side of its unique zero gives the conditional sorting
statement. The comparison with \(T/A\) is required because an \(I/E\) ranking
alone does not establish the globally selected route.

\paragraph{Sufficient conditions and zero-effect cases.}
The endpoint-crossing conditions are sufficient, not primitive conclusions.
If they fail, a finite interior cutoff need not exist. At low positive
\(k_i\), internal production remains feasible in principle but may be
arbitrarily unattractive because \(c_I\) and \(F_I\) are high. As
\(k_i\rightarrow\infty\), the maintained technology signs weakly favor
internal production relative to entrusted production. At \(M=0\), route \(E\)
is unavailable and there is no finite pre-MAH \(I/E\) cutoff.

\paragraph{Economic interpretation.}
Internal capability reduces both internal operating cost and setup burden.
Entrusted production is most useful to developers for which those internal
advantages are weak, but it is selected only when it also beats transfer or
delay. The result is organizational sorting, not a permanent firm-type
taxonomy.

\subsection{Proposition 2: The MAH-relevant project set}
\label{prop:p08-relevant-set}

\paragraph{Assumptions and fixed objects.}
Fix developer \(i\), project draw \((q,m)\), and support price \(p_m\).
The old route set is \(I,T,A\); post-MAH adds \(E\) without removing an old
option. Route values retain the definitions given above.

\paragraph{Definitions and claim.}
Define
\begin{equation}
  W_i^0(q,m)
  =
  \max\{W_i^I(q,m),W^T(q,m),0\},
  \qquad
  W_i^1(q,m;p_m)
  =
  \max\{W_i^0(q,m),W_i^E(q,m;1,p_m)\}.
  \label{eq:p08-old-new-values}
\end{equation}
Then
\begin{equation}
  W_i^1(q,m;p_m)-W_i^0(q,m)
  =
  \left[W_i^E(q,m;1,p_m)-W_i^0(q,m)\right]_+.
  \label{eq:p08-positive-part-gain}
\end{equation}
The MAH-relevant project set is
\begin{equation}
  \mathcal C_i(p_m)
  =
  \left\{(q,m):
  W_i^E(q,m;1,p_m)>W_i^0(q,m)\right\}.
  \label{eq:p08-relevant-set}
\end{equation}
Only project draws in this set receive a strict direct value gain at the fixed
support price.

\paragraph{Proof.}
For any real numbers \(a,b\),
\(\max\{a,b\}-a=[b-a]_+\). Apply the identity with
\(a=W_i^0\) and \(b=W_i^E\). The gain is strictly positive exactly where the
argument of the positive-part operator is positive, which is precisely
\(\mathcal C_i(p_m)\).

\paragraph{Sufficient conditions and zero-effect cases.}
A project receives a strict gain if and only if entrusted value strictly
exceeds internal value, transfer value, and abandonment value at the support
price. If \(W_i^E\leq W_i^0\), the project gain is zero. If this inequality
holds for all project draws, \(\mathcal C_i(p_m)\) is empty up to null sets.
As \(p_m\rightarrow\infty\), \(W_i^E\rightarrow-\infty\), so the old value is
preserved and the relevant set becomes empty.

\paragraph{Economic interpretation.}
MAH matters only where retained entrusted manufacturing beats every option
already available. The reform does not raise the value of projects that were
already better served internally, transferred, or delayed.

\subsection{Proposition 3: Project advancement and firm heterogeneity}
\label{prop:p08-advancement}

\paragraph{Assumptions and fixed objects.}
Fix a support price \(p_m\), policy-invariant primitives, and developer
characteristics except for the characteristic varied in each comparison.
Route values do not depend on research capability \(a_i\). For the
manufacturing-capability result, additionally impose the proposition-specific
sufficient condition \(\nu\geq1\).

\paragraph{Proof: expected value and advancement response.}
Let
\(\Omega_i^h(p_m)=\int W_i^h(q,m;p_m)\,dF(q,m)\) for
\(h\in\{0,1\}\), with the irrelevant \(p_m\) argument suppressed for
\(h=0\). Proposition~2 implies
\begin{equation}
  \Delta\Omega_i(p_m)
  \equiv
  \Omega_i^1(p_m)-\Omega_i^0
  =
  \int
  \left[W_i^E(q,m;1,p_m)-W_i^0(q,m)\right]_+
  dF(q,m)
  \geq0.
  \label{eq:p08-expected-gain}
\end{equation}
The common advancement response is
\begin{equation}
  \Delta x_i(p_m)
  =
  \left(\frac{\beta a_i}{\kappa}\right)^{1/\nu}
  \left[
    \left(\Omega_i^0+\Delta\Omega_i(p_m)\right)^{1/\nu}
    -
    \left(\Omega_i^0\right)^{1/\nu}
  \right].
  \label{eq:p08-advancement-change}
\end{equation}
Because \(t^{1/\nu}\) is strictly increasing for every \(\nu>0\),
\[
  \Delta x_i(p_m)>0
  \quad\Longleftrightarrow\quad
  \Delta\Omega_i(p_m)>0.
\]

\paragraph{Research-capability scaling.}
Holding \(k_i,p_m,\Omega_i^0,\Delta\Omega_i\) and all primitives fixed, a
positive response satisfies
\begin{equation}
  \frac{\partial\Delta x_i(p_m)}{\partial a_i}
  =
  \frac{\Delta x_i(p_m)}{\nu a_i}
  >0.
  \label{eq:p08-capability-response}
\end{equation}
Thus the level response is stronger at higher \(a_i\) whenever the
MAH-relevant set has positive expected surplus.

\paragraph{Capability result.}
The internal value is nondecreasing in \(k_i\), whereas the other old route
values and entrusted value are independent of \(k_i\). Therefore
\(\Omega_i^0(k_i)\) is nondecreasing and
\(\Delta\Omega_i(k_i;p_m)\) is nonincreasing. When \(\nu\geq1\),
\(h(t)=t^{1/\nu}\) is increasing and concave, so
\(h(u+d)-h(u)\) is nondecreasing in \(d\) and nonincreasing in \(u\).
Applying this result to \(u=\Omega_i^0(k_i)\) and
\(d=\Delta\Omega_i(k_i;p_m)\) yields
\begin{equation}
  k_i'\geq k_i
  \quad\Longrightarrow\quad
  \Delta x_i(k_i';p_m)
  \leq
  \Delta x_i(k_i;p_m),
  \qquad \nu\geq1.
  \label{eq:p08-manufacturing-response}
\end{equation}
The inequality is strict when the entrusted surplus falls on a
positive-measure set; for \(\nu>1\), a strict rise in positive baseline value
also makes the concavity channel strict. For \(0<\nu<1\), no sign is asserted
without an additional bound.

\paragraph{Zero-effect and boundary cases.}
If the relevant set is null or entrusted advantage is zero, then
\(\Delta\Omega_i=\Delta x_i=0\). At \(\nu=1\),
\[
  x_i^*=\frac{\beta a_i}{\kappa}\Omega_i,
  \qquad
  \Delta x_i=\frac{\beta a_i}{\kappa}\Delta\Omega_i,
\]
so manufacturing heterogeneity operates entirely through the expected value
gain. A high \(a_i\) does not create a response when the value gain is zero.

\paragraph{Economic interpretation.}
Research capability scales how strongly a given increase in the value of
advancing projects changes the level of advancement. Weak internal
manufacturing capability makes retained outsourcing more likely to add value.
The resulting response is strongest for high-\(a_i\), low- or
intermediate-\(k_i\) developers under the stated curvature condition. This
does not imply more patent applications, basic research, scientific
discovery, or breakthrough novelty.

\subsection{Corollary: Novelty composition is ambiguous}
\label{cor:p08-novelty}

\paragraph{Assumptions and fixed objects.}
Fix developer \(i\) and support price \(p_m\). The class
\(g\in\{O,\mathrm{Inc}\}\) is an exogenous classifier, not a control or state.
The same \(x_i\) applies to both classes.

\paragraph{Claim and proof.}
Define
\(\Omega_{ig}^0=E_{F_g}[W_i^0]\) and
\(\Omega_{ig}^1=E_{F_g}[\max\{W_i^0,W_i^E\}]\). Then
\begin{equation}
  \Delta\Omega_{ig}
  =
  E_{F_g}
  \left[
    \left(W_i^E-W_i^0\right)_+
  \right],
  \qquad
  \Delta\Omega_i
  =
  \sum_{g\in\{O,\mathrm{Inc}\}}\rho_g\Delta\Omega_{ig}.
  \label{eq:p08-class-gain}
\end{equation}
The positive-part identity proves the first equality; the mixture identity
for \(F\) proves the second. The baseline imposes no restriction ordering the
two conditional distributions of entrusted surplus. A support on which only
class \(O\) has positive entrusted surplus gives
\(\Delta\Omega_{iO}>\Delta\Omega_{i\mathrm{Inc}}\); interchanging the class
supports reverses the ranking. Hence neither ordering follows from the
baseline.

\paragraph{Zero-effect and boundary cases.}
If \(\rho_O=0\), class \(O\) contributes zero to the aggregate gain; if
\(\rho_{\mathrm{Inc}}=0\), class \(\mathrm{Inc}\) contributes zero. A
conditional class gain may be zero even when the other is positive.

\paragraph{Economic interpretation.}
Different novelty classes may have different manufacturing requirements,
commercial values, success probabilities, costs, or outside options. Only
explicit primitive differences can order their responses. No empirical
original-versus-incremental ranking is hard-coded into the theory.

\subsection{Proposition 4: CMO equilibrium existence and uniqueness}
\label{prop:p08-cmo-existence}

\paragraph{Assumptions and fixed objects.}
Fix \(M\) and all exogenous distributions and primitives. Use the maintained
market conditions: continuous weakly decreasing finite demand, continuous strictly
increasing supply, positive background demand at zero, vanishing demand at
high price, and unbounded high-price supply.

\paragraph{Claim.}
There exists a unique qualified-capacity clearing price
\(p_m^*(M)>0\).

\paragraph{Proof.}
Define excess demand
\begin{equation}
  Z_M(p_m)
  =
  D_m(p_m;M)-S_m(p_m).
  \label{eq:p08-excess-demand}
\end{equation}
It is continuous. At zero,
\(Z_M(0)\geq D_m^B(0)>0\). At sufficiently high prices, total demand tends to
zero and supply is unbounded, so \(Z_M(p_m)<0\). The intermediate value
theorem gives a positive root. Because demand is weakly decreasing and supply
strictly increasing, \(Z_M\) is strictly decreasing; the root is unique.

\paragraph{Sufficient conditions and zero-effect cases.}
Strictly increasing supply is sufficient for uniqueness even when demand has
flat regions. At \(M=0\), study-related entrusted demand is zero, but
background demand supports a positive pre-MAH price. If background demand at
zero were absent, the positive-price conclusion would not follow from these
conditions; that is outside the maintained baseline.

\paragraph{Economic interpretation.}
The capacity price coordinates deterministic entrusted-route demand,
project advancement, and supplier capacity through one scalar market. Its
existence does not require logit smoothing, supplier entry, or another market.

\subsection{Proposition 5: CMO scarcity attenuation}
\label{prop:p08-scarcity}

\paragraph{Assumptions and fixed objects.}
Let \(p_m^0=p_m^*(0)\) and \(p_m^1=p_m^*(1)\). Hold qualified-supply
technology, supplier distribution and background demand fixed across regimes.
The reform adds nonnegative entrusted-route study demand; old routes remain
available.

\paragraph{Proof: equilibrium price comparison.}
At the pre-MAH price,
\[
  Z_1(p_m^0)
  =
  D_m^{\mathrm{MAH}}(p_m^0;1)
  \geq0,
\]
because background demand equals supply at \(p_m^0\). Since \(Z_1\) is
strictly decreasing and vanishes at \(p_m^1\),
\begin{equation}
  p_m^1\geq p_m^0,
  \quad\text{with strict inequality if}\quad
  D_m^{\mathrm{MAH}}(p_m^0;1)>0.
  \label{eq:p08-price-order}
\end{equation}

\paragraph{Value attenuation.}
Define the fixed-price and equilibrium-price gains by
\[
  \Delta\Omega_i^{\mathrm{dir}}
  =
  \int
  [W_i^E(q,m;1,p_m^0)-W_i^0(q,m)]_+
  dF,
\]
\[
  \Delta\Omega_i^{\mathrm{eq}}
  =
  \int
  [W_i^E(q,m;1,p_m^1)-W_i^0(q,m)]_+
  dF.
\]
Since \(W_{i,p_m}^E=-b(m)<0\) and
\(p_m^1\geq p_m^0\), the integrand in the equilibrium expression is
pointwise no larger. Both integrands are nonnegative, so
\begin{equation}
  0
  \leq
  \Delta\Omega_i^{\mathrm{eq}}
  \leq
  \Delta\Omega_i^{\mathrm{dir}}.
  \label{eq:p08-scarcity-value-bounds}
\end{equation}

\paragraph{Advancement attenuation.}
For a fixed pre-MAH value \(\Omega_i^0\), the mapping
\[
  d\longmapsto
  \left(\frac{\beta a_i}{\kappa}\right)^{1/\nu}
  [(\Omega_i^0+d)^{1/\nu}-(\Omega_i^0)^{1/\nu}]
\]
is nonnegative and strictly increasing for \(d\geq0\). Therefore
\begin{equation}
  0
  \leq
  \Delta x_i^{\mathrm{eq}}
  \leq
  \Delta x_i^{\mathrm{dir}}.
  \label{eq:p08-scarcity-advancement-bounds}
\end{equation}

\paragraph{Zero-effect and boundary cases.}
If study demand is zero at \(p_m^0\), the price need not rise. If supply is
perfectly elastic at a common \(\bar p_m\), then
\(p_m^1=p_m^0=\bar p_m\) and attenuation disappears. If a high equilibrium
price eliminates entrusted advantage for developer \(i\), its equilibrium
gain is zero. Even then the old route-set value is preserved: scarcity can
attenuate the positive reform gain but cannot make access to the enlarged
choice set privately worse.

\paragraph{Economic interpretation.}
New demand for qualified external capacity bids up its scarcity price. This
reduces the value of the entrusted option relative to a fixed-price policy
calculation and tempers project advancement. The price response is endogenous;
MAH does not directly reduce scarcity or shift supply.

\subsection{Proposition 6: Planning-stage and observed outcomes}
\label{prop:p08-outcomes}

\paragraph{Assumptions and fixed objects.}
Fix a regime and its equilibrium objects. The common \(x_i^*\) is selected
before project characteristics and route assignment. The downstream
probabilities \(s(q)\) and \(s_g(q)\) are exogenous and policy invariant.
All integrals use the deterministic route choices evaluated in the same
regime.

\paragraph{Planning-stage projects.}
The equilibrium intensity at which viable projects reach the stage relevant
for commercialization is
\begin{equation}
  \Lambda_i^{\mathrm{plan}}
  =
  a_ix_i^*.
  \label{eq:p08-planning-arrival}
\end{equation}
This is the planning-stage technology
\(\lambda_i^{\mathrm{plan}}=a_ix_i\) evaluated at the advancement optimizer.

\paragraph{Observed retained outcomes.}
Within the paper's original-drug domain, expected realized outcomes for which
the developer remains the holder are
\begin{equation}
  Y_i^{\mathrm{ret}}
  =
  a_ix_i^*
  \int
  s(q)
  \left[
    \mathbf 1\{r_i^*=I\}
    +
    \mathbf 1\{r_i^*=E\}
  \right]
  dF(q,m).
  \label{eq:p08-retained-outcome}
\end{equation}
Expected realized outcomes with retained holder--producer separation are
\begin{equation}
  Y_i^E
  =
  a_ix_i^*
  \int
  s(q)\mathbf 1\{r_i^*=E\}
  dF(q,m).
  \label{eq:p08-entrusted-outcome}
\end{equation}
For \(g\in\{O,\mathrm{Inc}\}\), the class contribution to retained
outcomes is
\begin{equation}
  Y_{ig}^{\mathrm{ret}}
  =
  a_ix_i^*\rho_g
  \int
  s_g(q)
  \left[
    \mathbf 1\{r_i^*=I\}
    +
    \mathbf 1\{r_i^*=E\}
  \right]
  dF_g(q,m).
  \label{eq:p08-type-outcome}
\end{equation}
There is one common \(x_i^*\), not an \(x_{ig}\) control.

\paragraph{Proof and exact decomposition.}
Let
\[
  Q_i^{\mathrm{ret},h}
  =
  \int s(q)
  [\mathbf 1\{r_i^{*,h}=I\}+\mathbf 1\{r_i^{*,h}=E\}]
  dF
\]
and let \(x_i^h=x_i^*(h,p_m^h)\). Adding and subtracting
\(a_ix_i^1Q_i^{\mathrm{ret},0}\) gives
\begin{equation}
  \Delta Y_i^{\mathrm{ret}}
  =
  a_i(x_i^1-x_i^0)Q_i^{\mathrm{ret},0}
  +
  a_ix_i^1
  (Q_i^{\mathrm{ret},1}-Q_i^{\mathrm{ret},0}).
  \label{eq:p08-outcome-decomposition}
\end{equation}
The first term is the project-advancement/planning-arrival component. The
second is the retained-route composition component. Reassignment from \(I\)
to \(E\) can raise \(Y_i^E\) without changing the retained total; movement
from \(T/A\) to \(E\) can raise both.

\paragraph{Zero-effect and ambiguity cases.}
If advancement and all route indicators are unchanged, observed outcomes are
unchanged. A class-specific outcome can have zero change while the other
changes. The baseline allows
\(\Delta Y_{iO}^{\mathrm{ret}}=0<
\Delta Y_{i\mathrm{Inc}}^{\mathrm{ret}}\) and the reverse; it imposes no
ranking. If \(\rho_g=0\), that class contributes zero to the aggregate.

\paragraph{Economic interpretation.}
The model separates investment in advancing viable projects, their arrival at
the commercialization-relevant planning stage, organizational assignment, and
realized products. Holder--producer separation is observed only after route
assignment. None of these expressions makes patent generation endogenous or
lets MAH directly change downstream realization probabilities.

\subsection{Scope conclusion}
\label{subsec:p08-scope}

The six propositions and the novelty corollary exhaust the baseline results
established in this section. They add no market, no entry margin, no welfare object, no
logit route share, no class-specific advancement control, and no direct policy
effect on research capability, project draws, manufacturing technology,
capacity supply, or downstream realization.
